% Section: Technical Validation — Phase 2
% Author: Prachee Dhar (Task P2_04)
% IEEE Conference Paper Section - Card Fraud Triage (Phase 2)

\section{Technical Validation and Evaluation Protocol}
\label{sec:p2_validation}

This section reports the technical sanity checks, evaluation protocol details, and threshold policy for all Phase~2 models. The recorded model-selection criteria use validation statistics. Diagnostic sweeps also report test metrics for multiple configurations.

% ─── Evaluation Protocol ─────────────────────────────────────────────────────
\subsection{Evaluation Protocol and Anti-Leakage Controls}
\label{sec:p2_eval_protocol}

\paragraph{Chronological split.}
All models use the same chronological train/val/test split (by transaction time): train legit 170,524 records ($t \le 120{,}396$ s), val 56,961 records ($120{,}396 < t \le 145{,}247$ s, 57 frauds), test 56,962 records ($t > 145{,}248$ s, 75 frauds). No shuffling crosses split boundaries. Feature scaling, model fitting, and H3 representation scaling use legitimate training data. H2 normalizes scores using validation reference statistics.

\paragraph{Validation-only selection.}
Model hyperparameters (latent dim, $K$, covariance type, $\epsilon_{\text{cov}}$, fusion weight $\alpha$, noise level $\sigma$) are selected exclusively by validation PR-AUC. AE and DAE use validation-legitimate loss for early stopping and report the median validation-AP seed. Several sweep tables include test metrics for multiple configurations, so the record cannot establish a single sealed test evaluation.

\paragraph{Evaluation metrics.}
We report:
\begin{itemize}
  \item \textbf{PR-AUC}: Primary metric; threshold-free, appropriate for imbalanced detection (0.1\% base rate). 95\% bootstrap CI reported ($n=1{,}000$ resamples, seed=42).
  \item \textbf{ROC-AUC}: Included for cross-method comparability; can be misleadingly high under severe imbalance.
  \item \textbf{Recall@FPR}: Test ROC recall describes the ranking at a given test FPR. Thresholds fit on legitimate validation rows estimate operational workload, with separate realized test FPR and alert counts.
  \item \textbf{Precision@K}: Top-K precision at $K \in \{50, 100, 200\}$; directly measures triage quality for a fixed investigator capacity.
  \item \textbf{Paired bootstrap difference}: Tests whether PR-AUC improvement of a hybrid over its standalone parent is statistically significant ($n=1{,}000$, 95\% CI includes zero = not significant).
\end{itemize}

% ─── Seed Stability ──────────────────────────────────────────────────────────
\subsection{Seed Stability (5-Seed Experiment)}
\label{sec:p2_seed_stability}

Table~\ref{tab:p2_seed_stability} summarizes the five recorded runs. Across these seeds test AP varies substantially. The official AE and DAE files use the median seed ranked by validation AP.

\begin{table}[t]
\centering
\caption{5-Seed Stability: Val and Test PR-AUC (mean $\pm$ std).}
\label{tab:p2_seed_stability}
\footnotesize
\resizebox{\columnwidth}{!}{%
\begin{tabular}{lccc}
\toprule
\textbf{Model} & \textbf{Val PR-AUC} & \textbf{Test PR-AUC} & \textbf{n seeds} \\
\midrule
AE  & $0.269 \pm 0.132$ & $0.395 \pm 0.091$ & 5 \\
DAE & $0.279 \pm 0.169$ & $0.393 \pm 0.143$ & 5 \\
\bottomrule
\end{tabular}%
}
\end{table}

% ─── Gradient and Sanity Checks ──────────────────────────────────────────────
\subsection{Model Sanity Checks}
\label{sec:p2_sanity}

The following sanity checks confirm model integrity:

\begin{enumerate}
  \item \textbf{Gradient health:} No NaN gradients detected during AE/DAE training. Gradient norms converge to $[0.23, 0.81]$ range, indicating stable optimization without exploding or vanishing gradients.
  \item \textbf{Fraud/legit reconstruction gap:} The fraud-to-legit MSE ratio is 10.99$\times$ for the AE and 15.59$\times$ for the DAE on held-out batches, confirming the model reconstructs normal patterns well but fails on fraud patterns as intended.
  \item \textbf{Score direction:} All score files follow the convention: higher score = more anomalous. This was verified against the Phase~1 baseline for each model type.
  \item \textbf{PCA sanity check:} AE reconstruction error correlates with PCA reconstruction residual (Phase~1 PCA anomaly score), confirming the AE discovers similar structure to linear PCA while adding non-linear discrimination.
  \item \textbf{No test leakage:} Normalization statistics for DAGMM-lite's augmented representation $v$ use training legitimate records only. GMM parameters are fitted on training legit only. The ``val'' split label was used only for model selection, never for fitting parameters.
\end{enumerate}

% ─── Bootstrap CI ─────────────────────────────────────────────────────────────
\subsection{Paired Bootstrap Significance Testing}
\label{sec:p2_paired_bootstrap}

The paired bootstrap resamples transactions jointly across model scores. Against the stored ambient GMM, H2 improves AP by 0.2139 (95\% CI [0.1177, 0.3124]) and H3 by 0.1861 (95\% CI [0.0897, 0.2882]). The selected H2 is a monotonic transformation of the DAE score, so the H2 comparison describes the same ranking.

% ─── Runtime ──────────────────────────────────────────────────────────────────
\subsection{Runtime and Scalability}
\label{sec:p2_runtime}

\begin{table}[t]
\centering
\caption{Training and Inference Runtime (CPU, Apple Silicon). AE/DAE trained for 100 epochs; PCA/GMM/etc. include sweep time.}
\label{tab:p2_runtime}
\footnotesize
\resizebox{\columnwidth}{!}{%
\begin{tabular}{lrrl}
\toprule
\textbf{Model} & \textbf{Train (s)} & \textbf{Infer 10k (s)} & \textbf{Note} \\
\midrule
PCA Anomaly Score        & 0.44   & $<0.01$ & Full sweep \\
LOF                      & 4.42   & ---       & Sweep incl. \\
Isolation Forest         & 13.68  & ---       & Sweep incl. \\
OC-SVM                   & 58.08  & ---       & Sweep incl. \\
AE (n=3,531 params)      & 29.87  & 0.0018  & Best seed \\
DAE (n=2,895 params)     & 34.38  & 0.0014  & Best seed \\
DAGMM-lite (K=16, diag)  & $\sim$70 & 0.0021 & Grid search included \\
H2 Score Fusion          & $\approx0$ & $<0.001$ & Score combination only \\
\bottomrule
\end{tabular}%
}
\end{table}

% ─── Threshold Policy ─────────────────────────────────────────────────────────
\subsection{Operational Threshold Policy}
\label{sec:p2_threshold}

No operational threshold is recommended based on this research. We report threshold-free metrics (PR-AUC) and capacity-based metrics (Precision@K, Recall@FPR) to separate model quality from deployment decisions. Any threshold selection for production would require:
\begin{enumerate}
  \item Domain expert input on acceptable FPR (investigator workload) vs.\ recall (fraud recovery rate) trade-off.
  \item Re-calibration on a recent labeled window (concept drift monitoring).
  \item Separate cost/benefit analysis based on average transaction value and investigation cost.
\end{enumerate}

A cost curve analysis (varying the ratio of investigation cost to fraud loss) is available in the supplementary materials (\texttt{results/tables/eval\_cost\_curve.csv}), which confirms PCA Anomaly Score dominates across nearly all cost ratios examined.
